\documentclass[12pt,a4paper]{article}
\usepackage[russian]{babel}
\usepackage{fontspec}
\usepackage{geometry}
\usepackage{booktabs}
\usepackage{longtable}
\usepackage{array}
\usepackage{multirow}
\usepackage{enumitem}
\usepackage{titlesec}
\usepackage{xcolor}
\usepackage{siunitx}

\usepackage{pdfpages}

\usepackage{tablefootnote}

% Дополнительные пакеты для данного документа
\usepackage{caption}
\usepackage{subcaption}
\usepackage{listings}
% \usepackage{hyperref}

\geometry{a4paper, margin=25mm}
\setmainfont{Liberation Serif}
\setsansfont{Liberation Sans}

\definecolor{adblue}{RGB}{0,84,166}

\titleformat{\section}{\large\bfseries\color{adblue}}{}{0em}{}[\titlerule]
\titleformat{\subsection}{\normalsize\bfseries}{}{0em}{}
\titleformat{\subsubsection}{\normalsize\itshape}{}{0em}{}

\sisetup{output-decimal-marker={,}}

\usepackage{tikz}
\usetikzlibrary{positioning, calc, backgrounds, math, 3d, perspective, arrows.meta, graphs, graphdrawing, fit, shapes.geometric, trees}
    

\usepackage{pgfplots} % для построения графиков
\pgfplotsset{compat=1.18} % версия для совместимости
\usepackage{tkz-euclide}

% --- Подключение пакета алгоритмов (без флага algo2e, чтобы вернуть \begin{algorithm}) ---
\usepackage[ruled,vlined]{algorithm2e}
%\usepackage{caption}
\newcommand{\Gray}[1]{\textcolor{gray}{#1}} 
% Настройка подписей
\DeclareCaptionLabelFormat{gost}{#1 #2}
\captionsetup[table]{labelformat=gost,labelsep=endash,justification=justified,singlelinecheck=false,font=normal}
\captionsetup[figure]{labelformat=gost,labelsep=endash,justification=centering,font=normal}

% Локализация algorithm2e на русский (ОБЯЗАТЕЛЬНО после загрузки пакета)
\SetAlgorithmName{Алгоритм}{}{}
\SetAlgoCaptionSeparator{---}
\SetKwInput{Input}{Вход}
\SetKwInput{Output}{Выход}

\usepackage{url}
% Говорим пакету использовать моноширинный шрифт (как у texttt)
\Urlmuskip=0mu plus 1mu % гибкие пробелы для формул, если нужно
\DeclareUrlCommand\code{\urlstyle{tt}} 

% Библиотека схемотехники
\usepackage[siunitx, RPvoltages, american]{circuitikzgit}
\ctikzset{
    amplifiers/fill=cyan,
    sources/fill=green!20,
    diodes/fill=white,
    resistors/fill=violet,
    grounds/scale=1.2,
}

\usepackage{iftex}

\ifXeTeX % Если используется TeTeX или TeLaTeX
  \typeout{Режим XeTeX}
  \usepackage{xltxtra}
  \usepackage{fontspec}
  \usepackage{polyglossia}
  \setmainlanguage{russian}
  \setotherlanguage{english}
  \setkeys{russian}{babelshorthands=true} % По идее включает <<, >>, -- ---

  \setmainfont{Times New Roman}[Ligatures=TeX]
  \setromanfont{Times New Roman}[Ligatures=TeX]
  \setsansfont{Arial}[Ligatures=TeX]
  \setmonofont{Courier New}[Ligatures=TeX] 

  \newfontfamily{\cyrillicfont}{Times New Roman}
  \newfontfamily{\cyrillicfontrm}{Times New Roman}
  \newfontfamily{\cyrillicfonttt}{Courier New}
  \newfontfamily{\cyrillicfontsf}{Arial}
\fi

\ifLuaTeX % Если используется LuaLatex
  \typeout{Режим LuaTeX}
%  \usepackage{luatextra}
  \usepackage{fontspec}
  \defaultfontfeatures{Ligatures={TeX}}
  \usepackage{polyglossia}
  \setmainlanguage{russian}
  \setotherlanguage{english}
  \setmainfont{Times New Roman}[Script=Cyrillic, Ligatures=TeX]
  \setromanfont{Times New Roman}[Script=Cyrillic, Ligatures=TeX]
  \setsansfont{Arial}[Script=Cyrillic, Ligatures=TeX]
  \setmonofont{Courier New}[Ligatures=TeX]
  \newfontfamily{\cyrillicfont}{Times New Roman}[Script=Cyrillic, Ligatures=TeX]
  \newfontfamily{\cyrillicfontrm}{Times New Roman}[Script=Cyrillic, Ligatures=TeX]
  \newfontfamily{\cyrillicfonttt}{Courier New}
  \newfontfamily{\cyrillicfontsf}{Arial}[Script=Cyrillic, Ligatures=TeX]

  \usegdlibrary{trees}
  \usegdlibrary{force}
\fi

\ifPDFTeX % Если используется 8битный PDFTex
  \typeout{8-битный режим pdfTex}
  % For pdfTeX we must use old
  % 8-bit font technologies
  \usepackage[T2A]{fontenc}
  \usepackage[english,russian]{babel}
  \usepackage[utf8]{inputenc}
  \usepackage[english, russian]{babel}

  %\usepackage{newtxmath}
  \usepackage{amsmath}
  \usepackage{amsthm}
  \usepackage{amssymb}
  \usepackage{amsfonts}
  \usepackage{mathtext}
\fi

\usepackage{amsthm}  % Возможности AMSMath
\usepackage{amsmath}
\usepackage{amssymb}

\usepackage{esvect} % Для красивых стрелок
\usepackage{bm} % Для жирных символов

%Hyphenation rules
\usepackage{hyphenat}
\hyphenation{ма-те-ма-ти-ка вос-ста-нав-ли-вать ото-бра-жа-ют-ся сге-не-ри-ро-ван-ном сим-во-лы}
% \usepackage[english, russian]{babel}

\usepackage{graphicx}
\graphicspath{ {./images/} }

\usepackage{empheq}

\usepackage{float}

% \usepackage{comment}

\usepackage[hidelinks,
    unicode=true,          % поддерживает Unicode в закладках
    psdextra=true,      % автоматически конвертирует простую математику в текст
    % focusdefault=false
]{hyperref}


% Окружение "Теорема"
\newtheorem{theorem}{Теорема}[section]
\newtheorem{corollary}{Corollary}[theorem]
\newtheorem{lemma}[theorem]{Lemma}

% Окружение "Определение"
\theoremstyle{definition} % Прямой шрифт, нумерация
\newtheorem{definition}{Определение}[section]
\newtheorem{example}{Пример}[section]

% Окружение "Замечание"
\theoremstyle{remark} % Прямой шрифт, без нумерации
% Окуржение "Решение"
\newtheorem*{solution}{Решение}
\newtheorem*{remark}{Замечание}

\usepackage{polynom}
\usepackage{tabularx}
\usepackage{cancel}

\addto\captionsrussian{%
  \renewcommand{\figurename}{Рис.}%
  \renewcommand{\tablename}{Табл.}%
}

\renewcommand{\arraystretch}{1.5} % Высота строки таблицы

\renewcommand{\labelitemi}{\(\circ\)} % Бюллет -- пустой кружок в itemize


\begin{document}

\section{Задача}

\begin{itemize}
    \item Фотографии с большим количеством мелких деталей (кирпич,  Для обучения нужно 100 или 1000 картинок (используя поворот)  набор данных для обучения — резкие границы, на загубленных изображениях, где алгоритмы СМ МАТКАД!!! 
    \item Найти кодеки с описанием артефактов! 
    \item Анализ артефактов демозаики. 
    \item Сделать структуру — спектр структуры дискретизации. 
    \item Идея: описание как работает модель  пространственная (линза-массив светофильтров- матрица) 
    \item Матрица с дельта функциями в каждом отсчете. 
    \item Сделать такую структуру для у — по зелёному, и сделать структуру для Красного или синего. 
    \item Далее сделать спектр для всего этого для яркости и для цветности. 
    \item Далее свёртка.  В итоге получается, что некоторых частот нет, и мы достраиваем эти частотные диапазоны, например на базе обучения нейросети   
    \item Дообучить нейронку, сравнить уже имеющиеся датасеты, выбрать адекватный, и из них собрать что то свое, сохраняя структуру   
    \item Обработка
    \begin{itemize}
        \item Там где четная — зелёный, нужно сделать для красного и синего  
        \item Разработка обучающего датасета. 
        \item Фото RAW Берём картинку, уменьшаем хотя бы в 4 раза/8 раз разрешение в каждом направлении. 
        \item Добавляем шум (АБГШ) и без шум. 
        \item Делаем Аглументацию, и получаем много изображений, 
        \item После чего делаем демозаику и дальше уже сравниваем исходник с полученным изображением
        \item Какой шум добавить? Есть тестовое изображение — измеряем, получаем ЧКХ (получаем какую то характеристику), берём АБГШ, делаем свёртку для этой картинки
    \end{itemize}
\end{itemize}

\section{Теоретико-математическое моделирование физики сенсора и спектральный анализ процессов пространственной дискретизации}

\subsection{Сквозная пространственная физико-математическая модель формирования сигнала}
Процесс формирования цифрового изображения — это преобразование света в данные, описываемое цепочкой: «Объектив $\to$ Мозаика светофильтров (CFA) $\to$ Матрица фотоприемников». Объектив действует как пространственный фильтр нижних частот, описываемый функцией рассеяния точки ($\text{PSF}$), сворачивающей исходную сцену $f(x, y)$. Мозаика Байера селективно отсекает длины волн $\lambda$, редуцируя континуум в RGB-распределение, а матрица КМОП/ПЗС накапливает фотозаряды и квантует сигнал.

\subsection{Математическое описание дискретизации с использованием дельта-функций Дирака}
Для анализа в системах вроде MATLAB регулярная структура матрицы представляется 2D-решеткой дельта-функций Дирака. Идеальное сканирование моделируется как произведение непрерывной интенсивности $X(x, y)$ на этот модулирующий гребень, где каждый пиксель пропорционален интегральной яркости, упавшей на фотодиод.

\subsection{Математические структуры пространственной дискретизации для каналов яркости ($G$) и цветности ($R, B$)}
Байеровский фильтр создает пространственную маску. При шаге $h$, зеленый фильтр ($G$) использует шахматную структуру с функцией $M_G(x, y) = \sum_{m,n} \left[ \frac{1 + (-1)^{m+n}}{2} \right] \cdot \delta(x - mh, y - nh)$, обеспечивая высокую плотность отсчетов. Красные ($R$) и синие ($B$) фильтры занимают по 25\% сенсора, работая как разреженные структуры с шагом $h$ по осям.

\subsection{Спектральный анализ Фурье для яркостной и цветностной компонент дельта-структуры}
Применение преобразования Фурье ($\mathcal{F}$) переносит анализ в частотный домен $(u, v)$, где свертка порождает паразитные спектральные реплики. Спектр зеленого канала $S_G(u, v)$ содержит базовую низкочастотную составляющую и высокие частоты, смещенные на $1/(2h)$. Спектры $R$ и $B$ каналов подвержены сильному алиасингу из-за разреженности, образуя сложный набор накладывающихся интерференционных реплик.

\subsection{Анализ артефактов демозаики в терминах частотных диапазонов спектра}
Классическая интерполяция, действуя как ФНЧ, неизбежно вызывает артефакты из-за перекрытия спектров (aliasing). Основные дефекты — цветовой муар (ложные цвета от высоких частот), «молнии» (зазубрины на границах) и размытие деталей, возникающее при попытке подавить муар путем удаления высоких частот.

\subsection{Концепция восстановления отсутствующих частотных диапазонов на базе глубокого обучения нейросети NAFNet}
Информационная пустота в разреженных каналах устраняется сверточными нейросетями. Модель NAFNet, в отличие от линейных методов, обучается нелинейной многоканальной свёртке, восполняя частотные диапазоны через анализ пространственного контекста. Используя информацию канала $G$, сеть реконструирует $R$ и $B$ компоненты, а применение спектральных потерь (\textbf{Focal Frequency Loss}) позволяет восстанавливать высокие частоты, подавляя муар без потери резкости границ.

\section{Теоретико-математическое моделирование физики сенсора и спектральный анализ процессов пространственной дискретизации}

\subsection{Сквозная пространственная физико-математическая модель формирования сигнала}
Процесс формирования цифрового изображения в одноматричной видеосистеме представляет собой последовательное преобразование непрерывного светового поля в дискретный цифровой массив данных. Сквозная физическая модель описывается цепочкой элементов: \textbf{«Объектив (линза) $\to$ Массив пространственно-распределенных мозаичных светофильтров (CFA) $\to$ Полупроводниковая матрица фотоприемников»}.

\begin{enumerate}
    \item \textbf{Оптическая система (Линза)}: Объектив проецирует непрерывное спектральное распределение яркости сцены $f(x, y, \lambda)$ на плоскость сенсора. Оптическая система действует как пространственный фильтр нижних частот (ФНЧ). Ее интегральная характеристика описывается функцией рассеяния точки (ФРТ --- Point Spread Function, $\text{PSF}$). Математически изображение на входе в светофильтры $g(x,y)$ представляет собой пространственную свертку исходной сцены с ФРТ объектива:
    \begin{equation}
        g(x, y) = f(x, y) * \text{PSF}(x, y)
    \end{equation}
    
    \item \textbf{Массив светофильтров (CFA Байера)}: Световое поле проходит через мозаичную структуру, где в каждой пространственной координате $(x,y)$ отсекаются определенные длины волн $\lambda$. Континуум излучения редуцируется к дискретному пространственному распределению трех компонент: $R$ (красный), $G$ (зеленый) и $B$ (синий).
    
    \item \textbf{Матрица фотоприемников (КМОП/ПЗС)}: Осуществляет интегральное накопление фотозарядов на площади каждого пикселя и последующее квантование сигнала.
\end{enumerate}

\subsection{Математическое описание дискретизации с использованием дельта-функций Дирака}
Чтобы промоделировать этот процесс в системах компьютерной математики (SM Mathcad / MATLAB), регулярная структура идеальной матрицы представляется как двумерная решетка (гребень) $\delta$-функций Дирака. 

Пусть $X(x, y)$ --- непрерывное распределение интенсивности конкретного цветового канала в плоскости матрицы, а $\Delta x$ и $\Delta y$ --- пространственные шаги дискретизации (период расположения пикселей). Тогда математическая модель идеального пространственного отсчета (сканирования) запишется в виде:
\begin{equation}
    X_{discrete}(x, y) = X(x, y) \cdot \sum_{m=-\infty}^{\infty} \sum_{n=-\infty}^{\infty} \delta(x - m\Delta x, y - n\Delta y)
\end{equation}
Каждый отдельный пиксель в отсчете математически эквивалентен взвешенной дельта-функции, где вес --- это интегральная яркость, попавшая на фотодиод.

\subsection{Математические структуры пространственной дискретизации для каналов яркости ($G$) и цветности ($R, B$)}
Мозаика Байера накладывает пространственную маску на дельта-решетку. Опишем структуры пространственной дискретизации отдельно для каналов. Будем считать, что шаг дискретизации полноцветного изображения равен $h$.

\subsubsection{Канал яркости (Зеленый фильтр, $G$, Green)}
Зеленый фильтр занимает 50\% площади матрицы и расположен в шахматном порядке. Математическая структура модулирующей функции дискретизации $M_G(x,y)$ для зеленого канала имеет вид:
\begin{equation}
    M_G(x, y) = \sum_{m} \sum_{n} \left[ \frac{1 + (-1)^{m+n}}{2} \right] \cdot \delta(x - mh, y - nh)
\end{equation}
Четность суммы индексов $(m+n)$ обеспечивает шахматную структуру: фильтр пропускает сигнал только там, где $m$ и $n$ одновременно четные или одновременно нечетные.

\subsubsection{Каналы цветности (Красный $R$ и Синий $B$ фильтры)}
Красные и синие фотофильтры занимают по 25\% площади матрицы и располагаются в строках и столбцах через один отсчет.
Для красного канала $R$ (при допущении, что он находится в четных строках и четных столбцах):
\begin{equation}
    M_R(x, y) = \sum_{m} \sum_{n} \left[ \frac{1 + (-1)^m}{2} \cdot \frac{1 + (-1)^n}{2} \right] \cdot \delta(x - mh, y - nh)
\end{equation}
Для синего канала $B$ (смещен относительно красного на один шаг по вертикали и горизонтали):
\begin{equation}
    M_B(x, y) = \sum_{m} \sum_{n} \left[ \frac{1 - (-1)^m}{2} \cdot \frac{1 - (-1)^n}{2} \right] \cdot \delta(x - mh, y - nh)
\end{equation}

\subsection{Спектральный анализ Фурье для яркостной и цветностной компонент дельта-структуры}
Перейдем в частотный домен с помощью двумерного преобразования Фурье ($\mathcal{F}$). Прямое преобразование Фурье от модулирующих функций решетки порождает периодические спектральные реплики в частотной области $(u, v)$.

\subsubsection{Спектр яркостного канала ($\mathcal{F}\{M_G\}$)}
Поскольку зеленый канал имеет шахматную структуру, его спектр в частотной области содержит основную низкочастотную реплику в центре (постоянная составляющая и полезный сигнал яркости) и дополнительные высокочастотные реплики, смещенные на полупериоды частоты дискретизации. Математический спектр зеленого канала порождает модуляцию на граничной частоте Найквиста:
\begin{equation}
    S_G(u, v) = S_{base}(u, v) + S_{base}\left(u - \frac{1}{2h}, v - \frac{1}{2h}\right)
\end{equation}

\subsubsection{Спектр каналов цветности ($\mathcal{F}\{M_R\}, \mathcal{F}\{M_B\}$)}
Так как красный и синий каналы разрежены в 2 раза сильнее по осям $x$ и $y$, их частотные спектры содержат в 4 раза больше паразитных реплик (алиасинг). Их спектральные компоненты накладываются друг на друга на более низких пространственных частотах.

В итоге, полный спектр байеровского сигнала можно представить как суперпозицию базового спектра яркости (низкие частоты) и модулированных спектров цветности, смещенных в угловые и боковые частотные зоны.

\subsection{Анализ артефактов демозаики в терминах частотных диапазонов спектра}
Когда классический линейный алгоритм интерполяции (например, билинейный или бикубический) пытается восстановить полноцветный кадр, он выполняет операцию фильтрации нижних частот (сглаживание). Спектрально это означает наложение фиксированного окна ФНЧ.

Поскольку спектры яркости и цветности из-за дискретизации мозаикой Байера \textbf{частично перекрываются}, возникают фундаментальные артефакты демозаики:
\begin{enumerate}
    \item \textbf{Цветовой муар (Color Moire / Aliasing)}: Высокочастотная монохромная деталь (например, кирпичная стена или текстура ткани) из-за наложения спектров ошибочно интерпретируется алгоритмом как низкочастотная цветовая волна. На изображении появляются паразитные цветные полосы и пятна.
    \item \textbf{Артефакт «молнии» (Zipper effect)}: На резких границах объектов чередование пикселей Байера приводит к пространственному шатанию градиента. В частотной области это выражается в потере фазовой информации на предельных частотах Найквиста.
    \item \textbf{Размытие контуров (Blurring)}: Жесткая фильтрация спектра срезает все частоты выше частоты Найквиста разреженной решетки ($R/B$), стирая истинные резкие границы.
\end{enumerate}

\subsection{Концепция восстановления отсутствующих частотных диапазонов на базе глубокого обучения нейросети NAFNet}
Суть научной идеи и новизны исследования заключается в следующем: \textbf{информационная пустота в спектрах разреженных каналов $R, G, B$ может быть восполнена за счет аппроксимационных возможностей нейросети}.

Вместо математически глухого срезания частот фильтрами интерполяции, нейросеть (в нашем случае \textbf{NAFNet}) обучается нелинейной свёртке. Модель выступает в роли интеллектуального генеративного фильтра:
\begin{enumerate}
    \item Она анализирует пространственный контекст структуры мелких деталей (знает, как выглядит кирпич или геометрия стыков).
    \item На основе фазовых связей между выжившими отсчетами яркостногольного канала $G$, нейросеть производит \textbf{экстраполяцию и нелинейное достраивание частотных диапазонов} в каналах цветности ($R$ и $B$), которые были безвозвратно утеряны при дискретизации дельта-функциями.
    \item Оптимизация по частотному критерию (\textbf{Focal Frequency Loss}) заставляет скрытые слои сети штрафовать модель именно за несовпадение амплитудно-частотных характеристик (АЧХ) в области высоких частот, эффективно подавляя муар и восстанавливая резкость контуров, недостижимую для классических систем Mathcad.
\end{enumerate}

\section{Методология разработки синтетического обучающего датасета и математическое моделирование стохастических шумов сенсора}

\subsection{Обоснование структуры конвейера деградации исходных RAW-изображений}
Для формирования репрезентативного набора данных, обеспечивающего инвариантность нейросетевой модели к масштабу деталей и аппаратным искажениям, разработан специализированный конвейер деградации (downscaling/degradation pipeline). В качестве исходного материала ($GT$ --- Ground Truth) используются цифровые негативы высокого пространственного разрешения камеры Nikon D600.

Математический синтез парных изображений низкого качества ($LQ$) для обучения модели реставрации реализуется по следующему алгоритму:
\begin{enumerate}
    \item \textbf{Пространственное сжатие (Downscaling)}: Исходная высокодетализированная матрица $I_{HQ}$ изотропно уменьшается по каждому координатному направлению в $k$ раз (где $k \in \{4, 8\}$). Данная операция моделирует удаление малогабаритного сенсора от объекта съемки и искусственно переносит критически важные пространственные частоты сцены (контуры кирпичной кладки, мелкие текстуры) в дефицитную зону вблизи частоты Найквиста.
    \item \textbf{Частотно-модулированная фильтрация}: Перед интерполяцией выполняется свертка с ядром, имитирующим оптическое размытие и частотно-контрастную характеристику (ЧКХ) объектива.
    \item \textbf{Изолированное зашумление}: На полученный низкочастотный поток накладывается аддитивный белый гауссов шум, параметры которого привязаны к физическим характеристикам кремниевого сенсора.
\end{enumerate}

\subsection{Математическое моделирование аддитивного белого гауссова шума на основе ЧКХ}
Для генерации физически достоверного шума в рамках НИР применяется подход, связывающий дисперсию флуктуаций с частотно-контрастной характеристикой реальной оптической системы. Реальный цифровой шум на изображении не является чисто белым (попиксельно независимым), так как он проходит через передаточную функцию тракта считывания.

Математическая модель стохастического шума $n(x, y)$ базируется на измерении передаточной функции модуляции (MTF) по тестовой мире. Пусть $H_{MTF}(u, v)$ --- двухмерный профиль ЧКХ системы. Тогда генерация шума описывается как пространственная свертка некоррелированного АБГШ $\xi(x, y)$ с импульсной характеристикой оптико-электронного тракта $h_{CFA}(x, y)$:
\begin{equation}
    n(x, y) = \xi(x, y) * h_{CFA}(x, y)
\end{equation}
где $\xi(x, y) \sim \mathcal{N}(0, \sigma^2)$, а профиль функции $h_{CFA}(x, y)$ определяется через обратное преобразование Фурье от ЧКХ:
\begin{equation}
    h_{CFA}(x, y) = \mathcal{F}^{-1}\{H_{MTF}(u, v)\}
\end{equation}

Итоговое синтезированное изображение низкого качества с учетом ЧКХ-модулированного шума принимает вид:
\begin{equation}
    I_{LQ}(x, y) = \text{Downsample}_k \Big( I_{HQ}(x, y) * h_{lens}(x, y) \Big) + n(x, y)
\end{equation}

\subsection{Пайплайн многоканальной аугментации и генерации тренировочных патчей}
Полнокадровые матрицы разрешением 24 Мп не могут быть обработаны нейросетью NAFNet монолитно из-за ограничений локальной видеопамяти графического процессора (GPU OOM). В модуле \texttt{prepare\_dataset.py} реализован алгоритм дискретной нарезки на перекрывающиеся суб-изображения (патчи) фиксированного размера $256 \times 256$ пикселей.

Для предотвращения эффекта переобучения и кратного искусственного увеличения объема выборки (до $10^3$---$10^4$ эквивалентных кадров) применяется комбинация геометрических трансформаций из библиотеки \texttt{Albumentations}. Пайплайн включает:
\begin{itemize}
    \item \textbf{Синхронные пространственные инварианты}: Горизонтальные (\texttt{HorizontalFlip}) и вертикальные (\texttt{VerticalFlip}) отражения, а также повороты строго на углы, кратные 90 градусам (\texttt{RandomRotate90}). Применение строго ортогональных поворотов гарантирует сохранение фазовой структуры пикселей мозаики Байера и исключает артефакты интерполяции краев.
    \item \textbf{Асинхронные цветовые трансформации}: Стохастическое изменение яркости и контраста, применяемое исключительно к LQ-потоку для имитации нестабильности экспонометрической системы малогабаритного сенсора.
\end{itemize}

\subsection{Сравнительный анализ и интеграция существующих наборов данных}
В рамках исследования проведен аналитический обзор открытых академических датасетов (DIV2K, Flickr2K, SIDD), ориентированных на задачи Super-Resolution и Denoising. Выявлено, что стандартные наборы данных не учитывают специфику декомпозиции несжатых RAW-данных камеры Nikon D600 во временном и пространственном доменах.

Научно обосновано решение по формированию гибридного датасета:
\begin{enumerate}
    \item Из верифицированных внешних датасетов извлекаются эталонные высококонтрастные текстурные маски ($GT$).
    \item Геометрическая структура и спектральный профиль шума полностью пересчитываются программным комплексом на основе кастомных параметров ЧКХ сенсора Nikon D600, полученных в ходе эксперимента.
    \item Финальный массив данных сохраняется в виде бинарных баз данных \texttt{LMDB}, что обеспечивает линейную скорость чтения тензоров загрузчиком \texttt{DataLoader} PyTorch во время итераций оптимизации.
\end{enumerate}

\section{Архитектура нейросетевого вычислительного ядра NAFNet и оптимизация критериев качества реставрации}

\subsection{Обоснование выбора и математическое описание блоков NAFNet}
Для сквозной пространственно-цветовой реставрации изображений в рамках данного исследования интегрирована сверточная нейросетевая архитектура \textbf{NAFNet} (Nonlinear Activation Free Network). Ее ключевое теоретическое отличие от классических архитектур (например, RCAN, SwinIR) заключается в полном отказе от традиционных нелинейных функций активации (ReLU, GELU) в пользу вычислительно эффективных алгебраических операций над тензорами признаков. Это позволяет избежать проблемы затухания градиентов и потери тонких текстурных переходов при глубокой свертке RAW-данных.

Магистральный блок NAFNet базируется на концепции упрощенного вентильного механизма (Simplified Channel Attention) и нелинейного перемножения признаков. Пусть $X \in \mathbb{R}^{H \times W \times C}$ --- входной тензор признаков на определенном слое сети. Математическая логика работы нелинейного блока \textbf{Gated Linear Unit (GLU)} без функции активации описывается следующим образом:
\begin{enumerate}
    \item Тензор признаков проецируется через линейный слой (свертку $1 \times 1$) для удвоения канальной размерности: $X' = \text{Conv}_{1\times1}(X) \in \mathbb{R}^{H \times W \times 2C}$.
    \item Полученный массив разделяется по оси каналов на два равных подматричных тензора: $X_1, X_2 \in \mathbb{R}^{H \times W \times C}$.
    \item Нелинейность реализуется через операцию поточечного произведения (перемножения Адамара):
    \begin{equation}
        \text{Gate}(X) = X_1 \odot X_2
    \end{equation}
\end{enumerate}

Для учета глобального контекста и взаимного влияния цветовых каналов $R, G, B$ в блоке NAFNet применяется упрощенный механизм внимания к каналам (\textbf{Simplified Channel Attention --- SCA}). Вместо сложной двухслойной структуры MLP, вычисление весов каналов производится через операцию глобального усреднения (Global Average Pooling, GAP) и одну линейную проекцию:
\begin{equation}
    \text{SCA}(X) = X \cdot \text{Conv}_{1\times1}(\text{GAP}(X))
\end{equation}

Такая структура минимизирует вычислительную сложность алгоритма, позволяя эффективно восстанавливать высокочастотные диапазоны спектра в режиме реального времени.

\subsection{Математическая модель комбинированного критерия оптимизации}
В ходе НИР теоретически обосновано и практически доказано, что использование исключительно пространственных функций потерь ($L_1$/MSE) приводит к размытию контуров. Для восстановления резких границ (например, стыков кирпичной кладки, мелкого цифрового зерна) в скрипте \texttt{smoke\_test.py} и цикле обучения \texttt{run\_training.py} реализована комбинированная функция потерь.

Полный функционал потерь представляет собой взвешенную линейную комбинацию пространственного и частотного критериев:
\begin{equation}
    \mathcal{L}_{total} = \mathcal{L}_{L1} + \gamma \mathcal{L}_{FFL}
\end{equation}
где $\gamma$ --- масштабирующий гиперпараметр, регулирующий вклад частотной компоненты.

\subsubsection{Пространственная компонента ($\mathcal{L}_{L1}$)}
Обеспечивает базовое совпадение геометрии и глобальной контрастности на уровне пикселей:
\begin{equation}
    \mathcal{L}_{L1} = \frac{1}{H \times W \times C}\sum_{x,y,c} \left| I_{HQ}(x,y,c) - I_{pred}(x,y,c) \right|
\end{equation}

\subsubsection{Частотная компонента ($\mathcal{L}_{FFL}$ --- Focal Frequency Loss)}
Подавляет цветовой муар и восстанавливает утерянные при дискретизации частоты. Алгоритм вычисляет двумерное дискретное преобразование Фурье ($\text{2D-DFT}$) от эталона и предсказания сети:
\begin{equation}
    F(u,v) = \mathcal{F}\{I(x,y)\} = \sum_{x=0}^{H-1}\sum_{y=0}^{W-1} I(x,y) \exp\left( -j2\pi \left( \frac{ux}{H} + \frac{vy}{W} \right) \right)
\end{equation}
На основе комплексных спектров $F_{HQ}(u,v)$ и $F_{pred}(u,v)$ формируется динамическая матрица пространственно-частотной фокусировки $w(u,v)$, которая увеличивает штраф за ошибки на тех частотах, где расхождение спектров максимально (высокочастотные шумы, резкие градиенты):
\begin{equation}
    w(u,v) = \left| F_{HQ}(u,v) - F_{pred}(u,v) \right|^\alpha
\end{equation}
Математическое выражение для $\mathcal{L}_{FFL}$ принимает вид:
\begin{equation}
    \mathcal{L}_{FFL} = \frac{1}{H \times W}\sum_{u=0}^{H-1}\sum_{v=0}^{W-1} w(u,v) \left\| F_{HQ}(u,v) - F_{pred}(u,v) \right\|^2
\end{equation}

\subsection{Методика дообучения (Fine-tuning) и трансферного переноса знаний}
Для достижения максимального качества реставрации без чрезмерных временных затрат на расчеты применяется двухэтапная стратегия обучения (Transfer Learning):
\begin{enumerate}
    \item \textbf{Предобучение (Pre-training)}: Модель NAFNet инициализируется базовыми весами, полученными в ходе обучения на крупномасштабных открытых наборах данных (например, DIV2K). На этом этапе сеть осваивает общие пространственные закономерности естественных сцен.
    \item \textbf{Адаптивное дообучение (Fine-tuning)}: На втором этапе веса модели настраиваются строго на кастомном подмножестве патчей, полученных с камеры Nikon D600 через скрипт \texttt{prepare\_dataset.py}. Обучение регулируется планировщиком \texttt{TrueCosineAnnealingLR}, динамически снижающим шаг обучения до $\eta_{min}=10^{-7}$ на протяжении $400\,000$ итераций, что предотвращает разрушение предобученных признаков (catastrophic forgetting).
\end{enumerate}

\subsection{Метрики объективной оценки качества реставрации}
Экспериментальная оценка разработанных методов в модуле \texttt{evaluate\_baseline.py} базируется на двух фундаментальных математических метриках качества изображений:


\begin{enumerate}
    \item \textbf{Пиковое отношение сигнал/шум (PSNR --- Peak Signal-to-Noise Ratio)}: Оценивает попиксельную точность восстановления амплитуды сигнала и вычисляется через среднеквадратичную ошибку ($\text{MSE}$):
    \begin{equation}
        \text{PSNR} = 10 \cdot \log_{10} \left( \frac{\text{MAX}_I^2}{\text{MSE}} \right)
    \end{equation}
    где $\text{MAX}_I = 1.0$ --- максимальное теоретическое значение нормированного пикселя тензора.
    
    \item \textbf{Индекс структурного сходства (SSIM --- Structural Similarity Index)}: Моделирует особенности человеческого зрительного восприятия, оценивая искажения трех независимых компонент кадра --- яркости ($l$), контраста ($c$) и структуры ($s$):
    \begin{equation}
        \text{SSIM}(x, y) = \frac{(2\mu_x\mu_y + C_1)(2\sigma_{xy} + C_2)}{(\mu_x^2 + \mu_y^2 + C_1)(\sigma_x^2 + \sigma_y^2 + C_2)}
    \end{equation}
    где $\mu$, $\sigma^2$ и $\sigma_{xy}$ --- средние значения, дисперсии и ковариация локальных областей сравниваемых изображений, а $C_1, C_2$ --- стабилизирующие константы.
\end{enumerate}


\end{document}